\section{Sex and Liberation}

\begin{quotex}
It is no secret now that lust is also a form of addiction. My point here is that the current regime knows this and exploits this situation to its own advantage. In other words, sexual freedom is really a form of social control….The truth, the one propagated by the regime through advertising, sex education, Hollywood films, and the university system — the truth, in other words, for general consumption — is that sexual liberation is freedom. The esoteric truth, the one that informs the operations manual of the regime — in other words the people who benefit from ``liberty" — is the exact opposite, namely, that sexual liberation is a form of control, a way of maintaining the regime in power by exploiting the passions of the naïve, who identify with their passions as if they were their own and identify with the regime which ostensibly enables them to gratify these passions. \flright{\textsc{E Michael Jones}, \emph{Libido Dominandi}\footnote{\url{http://culturewars.com/CultureWars/1999/libido_ad.html}}}

These began with divine ideas by way of contemplation of the heavens with the bodily eyes. Thus in their science of augury the Romans used the verb \emph{contemplari} for observing the parts of the sky whence the auguries came or the auspices were taken. \flright{\textsc{Giambattista Vico}, \emph{The New Science}}

\end{quotex}
\paragraph{Intro}
During my youth, I used to study the writings of Herbert Marcuse, sometimes called the ``Julius Evola of the Left." There is some justice in that appellation since their interests overlapped — fascism, sexuality, the drug culture, hippies, etc. — and they both opposed the modern world, particularly in the dual pincers of American capitalism and Soviet communism.

In \emph{Reason and Revolution}, his study of Hegel, Marcuse opposed the monopoly of rationalism by positive science. Thus philosophy or metaphysics has a claim on truth and rationality. That has stuck with me, particularly because my fields of study were maths, physics, and chemistry.

In \emph{One Dimensional Man}, he documented the deleterious effects of consumer society on human fulfillment. I was especially interested in \emph{Eros and Civilization}. Marcuse combined the personal and the social in a powerful way. That technique provides a lesson for the development of a historiography of the right. For example, we've related Plato's political systems to the rule of men dominated by the different soul centers.

This brief essay continues that project, albeit in a sketched-out form. For the reconstruction of the prehistory of the West, I am relying on \textbf{Giambattisa Vico}'s \emph{New Science} as well as \textbf{Fustel de Coulange}'s \emph{Ancient City}, as well as the studies of myths by \textbf{J J Bachofen}, \textbf{Julius Evola} and \textbf{Julian Jaynes}.

\paragraph{Eros and Civilization}
At that time, it was possible to be a materialist atheist using Darwin, Freud and Marx as the foundation. Life, mind, and society could be explained without recourse to anything transcendent. Marcuse combined Freud and Marx very effectively. Marx had always claimed that consciousness was formed by the material conditions of life, without explaining how. Marcuse took the major themes of Freudian psychology to show how social structures become reflected in human psychology.

Thus Freud's major themes of sexuality and repression could be explained in the individual by the systems of production and consumption. Since Freud's conceptions are ultimately conservative, Marcuse had to rework his ideas quite a bit.

As is common with the left, Marcuse felt that men are born good but become repressed by society. Unlike Freud who contended that the sexual instinct (eros) had to be sublimated for civilization to develop, Marcuse claimed that unleashed sexuality is liberating and constructive. For Freud, the superego, by absorbing cultural norms, was able to repress sexuality. However, for Marcuse, the superego instead absorbed unjust social structures. This ``surplus repression" was beyond anything really necessary.

The Freudian system apes the traditional view of man. The ego is the self, or rational center of consciousness. The superego is what is transcendent to the ego and the id represents desires and instincts. Of course, for Freud there is no real transcendence, hence it can be explained only by material terms; that provides the opening for Marcuse.

So if the superego is merely the result of alienated consciousness, there is no reason at all to suppress the instincts. The assumption must be that they are good in themselves, which really means that Darwinian random evolution somehow produces the good.

Freud said that we go through three stages: oral, anal, genital. In maturity, our sexuality is centered on the genital phase. Marcuse regarded that as a form of repression. Hence, in his utopian vision of human centered production based on technology, the fully developed human would become polymorphous perverse, like a child. His whole body would be sexually charged, without the repression of the oral and anal stage.

That hardly seems like a manly or worthwhile goal, but there is little doubt that such an idea resonates in contemporary culture. This is evident from the porno film \emph{Deep Throat} in the 1970s which mainstreamed fellatio to more recent obsession with the ass and anal sex. A recent episode of the ``edgy" HBO series \emph{Girls} featured a gratuitous scene of anilingus.

This is an inversion of true Eros which is a drive to God and the transcendent. Biological sex, then, is a mere distortion and materialization of that drive.

\paragraph{Fundamental Principles}
Giambattista Vico identifies three customs of all nations:

\begin{itemize}
\item All have some religion 
\item All contract solemn marriages 
\item All bury the dead 
\end{itemize}
These, therefore, are the basis of the first three principles of his New Science. He points out that they are reinforced by their religious practices. The consequences of not doing so are, according to Vico:

\begin{quotex}
Otherwise, the world would return to a brutish state and again become a wilderness. 

\end{quotex}
Now the facts of human societies can be reformulated as principles:

\begin{itemize}
\item Divine Providence 
\item The moderation of passions through marriage 
\item The immortality of human souls attested by burial 
\end{itemize}
Just to be clear, Vico again emphasizes:

\begin{quotex}
These are the boundaries of human reason and transgressing them means abandoning our humanity. 

\end{quotex}
Of course, modernity in its fullness is the exact opposite:

\begin{itemize}
\item Man to be fully human requires freedom from God 
\item The free expression of passions, especially but not only sexuality, is human liberation 
\item Live for the moment 
\end{itemize}
Thus, we have another barometer to measure the decline, or progress if you prefer, of the human race. Of course, most people do not consciously think about such things and try to steer a middle, i.e., a ``lukewarm", course; but those are the positions for those who think it through.

\paragraph{Providence and Religion}
Vico makes Providence a principle of his New Science, so he regards Providence as a historically veribiable fact. He writes:

\begin{quotex}
Providence must provide a history of the orders and institutions which providence bestowed on the great polity of humankind without the knowledge or advice of humankind, and often contrary to human planning. For although by its creation our world is temporal and particular, the orders which providence establishes in it are universal and eternal. 

\end{quotex}
Thus he claims that the human race could not have survived without divine providence. The actions of men are insufficient to explain all the cultural and social benefits that accrued to them.

The modernist, on the other hand, rejects Providence. For him, the human race evolved through random processes into an environment that quite fortuitously was able to support human life. Gods and religions, therefore, are social constructs used to bolster and legitimize existing social arrangements.

\textbf{Rene Guenon} expressed it this way:

\begin{quotex}
It is not that the order of the cosmos was conceived on the model of social institutions, but, on the contrary, these institutions were founded on the basis of analogy with the cosmic order. 

\end{quotex}
\paragraph{Marriage and Immortality}
\begin{quotex}
When we discuss immortality, we must grant considerable importance to the consensus of humankind, who either fear or worship the spirits of the underworld. I follow this general belief. \flright{\textsc{Seneca}}

\end{quotex}
Vico recognized a hierarchy of allegiances. First a man cared for himself, then his wife and children. This extended to his clan, nation and so on. Marriages were solemnized not least of all for the sake of the children. Without that legal bond, the children could be left to their own devices. Vico was prescient in foreseeing that the children of unwed mothers would come to depend on public or private charity, as we see today. Any exceptions just prove the rule.

Vico simply assumes that sexual relationships are genital. In the Freudian scheme, the child advances sexually in three phases: oral, anal, and genital. Hence, the psychologically healthy and rational adult regards sex as genital-to-genital contact. Its proper use is within marriage as he describes. Obviously, this is not an arbitrary command of a totalitarian church, but rather the common experience of the human race, insofar as it is civilized and guided by divine providence. Specifically, this applies also to pagan societies pace the new age neo-pagans, and Evola alludes to it in Ancient Rome:

\begin{quotex}
every offense to the sacredness of aristocratic marriage and to the lineage was considered as a crime above all in the face of the \emph{genius} of the lineage 

\end{quotex}
On the contrary, the revolutionary goal is to return to the childlike innocence of polymorphous perversity, for whom the entire body is sexually charged. This is considered liberation although, as E Michael Jones points out, in its actual effect it is a form of social control. Rather than the result of the overthrowing of capitalism as Marcuse had hoped, sexual control becomes just another capitalist commodity. In Marcuse's terms, sexual liberation has been ``co-opted" by the capitalist system he opposed.

\paragraph{Historical Reconstruction}
So, in his reconstruction of prehistory, Vico sees the formation of society the father, as the protector of his family. The fathers gather together for mutual support and protection. Thus, at its origin, society was a patriarchy. This is not exactly the mannerbund (with is homoerotic overtones in some circles), as Evola would have it, since it was composed strictly of fathers. Evola does recognize the ``son of duty", since that is the only way that the father's memory would continue in the burial rites.

Now each family was considered to have a divine origin as the primal father was regarded as a god. The consequence is that each family also had its own religion with rites and legends of its own. The clan also had a divine origin and its father, not necessarily biological but certainly spiritual, was likewise considered a god. The result was a hierarchy of religions and religious practices. Daily life was regulated by the rites, taboos, and rituals of the religions.

Thus, society was at its origin patriarchal and theocratic. Vico called it the age of the gods. Clients were attracted to the patriarchy for protection. They did not have the same spiritual foundation. Hence, they followed the rites of whatever family they are attached to. They were not necessarily biologically related; over time, due to endogamous marriage, racial differences between the clients and the families become more pronounced.

They didn't have an inner center (ego in Freud's terms) and were driven primarily by desire, or id, in Freudian sense. They may even have been pre-Adamics. Their consciousness was formed by the families. Without a strong ego, or rational capability, they are directed by command, poetry, symbolism (see Jaynes). This inner consciousness of external commands is the Freudian superego.

Over time, the age of the gods gave way to the age of the heroes. This was the result of a schism between the priestly and noble functions. In Platonic terms, it was the movement from aristocracy to timocracy. The next phase was toward democracy or popular rule. This was the result of the desire of the clients and plebeians for more political and social control.

\paragraph{Natural and Conventional Relationships}
Oddly enough, the revolutionary and reactionary reconstructions of history are not that different, although their evaluations certainly are. This is not totally surprising, since Marx himself relied on Vico's works. The revolutionary sees the overturning of patriarchy and theocracy as forward progress. Each stage of decline for the reactionary is progress for the revolutionary. More and more people take part in social life.

For the reactionary, the stages move downward from the rational (in the higher sense) to the emotional to desire and sensation. The latter has three substages: normal desires, excessive desires, and unnatural desires. So the aristocrat knows the divine order of the cosmos and orders social relationships in conformance to it. The democrat lacks that understanding. Originally, aristocratic knowledge came to the democrat through poetry, symbols, and divine commands. In his consciousness, it was experienced as alien and alienating, since it did not comport with the other parts of his interiority. In order to overthrow the commands of the superego, he also had to eliminate its source in religion and patriarchy.

For the reactionary, theocracy, patriarchy and hierarchy are natural and organic structures, a reflection of the divine order. For the revolutionary, they are merely conventional and social structures. They have no origin beyond the human mind. At worst, they serve to maintain existing power structures. For Marx, they are the product of economic arrangements. In a state of wild nature, a patriarchy may have made some sense. However, as economic activity became more complex, power in the hands of the few was not justifiable.

To close the circle with Marcuse: the Old Left was failing because the prosperity in the West of the working classes tended to make them more conservative rather than revolutionary. Current arrangements suited them, so they could not be easily motivated to revolt. Hence, the New Left looked for the revolutionary consciousness not in the working class, but in other groups considered to be marginalized by social arrangements.

These thus are the choices: either the arrangements are rational and supported by divine providence or they are irrational and arbitrary. There can be no middle ground.



\flrightit{Posted on 2015-01-22 by Cologero }

\begin{center}* * *\end{center}

\begin{footnotesize}\begin{sffamily}



\texttt{White Rabbit on 2015-01-23 at 01:05 said: }

Sex as a form of mind control? I don't know. I've considered it, but when I talk about it my friends tell me I need to get laid.


\hfill

\texttt{X on 2015-01-23 at 02:33 said: }

Good article.

Actually the result of sexual liberation is the dissolution of true sexuality, everything which rejects being will finally fall into absolute non-being. The more sexuality is liberated, the less polarity exists between sexes, the less true sexuality there is.


\hfill

\texttt{David on 2015-01-23 at 04:32 said: }

Interesting articles with differents things to think upon, especially the study of ancient family regarding rites. 

But I have one question : how does the "Freudian system apes the traditional view of man"? It is my impression that the Freudian system is a reversal and misunderstanding of the traditional view of man since for Freud, the conscious (let alone any thought of supra-conscious or supra-individual !) actions are dictated by the drive of the subconscious and its energy. If I do action X, it is because I have that amount of energy which made my reprensentation X bypass my "screen". There is no really conscious choice or action for Freud that seems possible if not for the creation of said "screens" of the pulsion and passion from a "super-ego" which in the end is no more than a construction and by no mean a Self in the higher sense. At the same time, I am not a student of Freud, and I would therefore ask for more guidance on this.


\hfill

\texttt{Cologero on 2015-01-23 at 05:09 said: }

In the same sense that `Satan is the ape of God.' See Guenon's Reign of Quantity:

\begin{quotex}
by virtue of the law of analogy, the lowest point is the obscure reflection or the inverted image of the highest point, from which follows the consequence, paradoxical only in appearance, that the most complete absence of all principle implies a sort of counterfeit of the principle itself, something that has been expressed in a theological form in the words `Satan is the ape of God'.

A proper appreciation of this fact can help greatly toward the understanding of some of the darkest enigmas of the modern world, enigmas which that world itself denies because, though it carries them in itself, it is incapable of perceiving them, and because this denial is an indispensable conditions for the maintenance of the special mentality whereby it exists.

If our contemporaries as a whole could see what it is that is guiding them and where they are really going, the modern world would at once cease to exist as such, for the rectification could not fail to come about through that very circumstance. 

\end{quotex}
That is central to our project.


\hfill

\texttt{David on 2015-01-23 at 08:44 said: }

@Colgero : Then I am sorry to have misunderstood your purpose. I must say that the word "apes" was unknown to me, so I tried to understand the sentence without it, which led me to be believe it had to do with "copying"; hence my answer asking why you linked Tradition to the Freudian project. Thank you for the answer.


\hfill

\texttt{Cologero on 2015-01-23 at 09:13 said: }

The word was chosen as an allusion to Guenon's usage.

In French: « Satan est le singe de Dieu »

Rene Guenon \textit{Le Regne de la Quantite et les Signes des Temps}

\url{http://ekladata.com/ZvjZowigoi2MzLb65eOL92PGSD0/Rene-Guenon-1945-Le-Regne-de-la-Quantite-et-les-Signes-des-Temps.pdf}


\hfill

\texttt{Max on 2015-01-23 at 12:12 said: }

Cologero, you mention clients ruled by desire and contrast them to the patriarchal families. How are we to understand that the Roman families traced their ancestry to Aeneas who decended from Venus (the word related to the sanskrit ``vanas", desire, and to the scandinavian Vanir who were the mythological ancestors of the Yngling dynasty of Swedish kings). Perhaps as a conscious memory of primordiality and virginal nature since the world is apart from desire more significantly perpetuated by love.

Here there is relevance to your quote of Guenon about the highest reflected in the lowest as both the patricians and plebeians are apparently ``descended" from desire, but understood in different ways. The lower kind is in truth just a substitute, a realization which opens a possibility to overcome the denial. This suggests that the common man may recover his nobility through a purification of desire and earn the royal privilege to claim descent from Our Lady and the Heavenly Father.


\hfill

\texttt{Cologero on 2015-01-23 at 14:33 said: }

Max, I'll give you Vico's explantion [512]:

\begin{quotex}
Venus was also regarded as a goddess of solemn marriages. Called pronuba, bridesmaid, she covered her private parts with her famed girdle, which lascivious poets later embroidered with all sorts of aphrodisiac charms. At that point, the severe historical truth of the auspices had been corrupted, and people believed that Venus lay with men, as Jupiter lay with women. Hence, Aeneas, who was conceived under Venus' auspices, was said to be Venus' son by Anchises. 

\end{quotex}
So being ``conceived under Venus' auspices" is different from what you suggest. The poets represent a corruption.


\hfill

\texttt{Tom Blanchard on 2015-01-23 at 15:03 said: }

With regard to the burial of the dead as the hallmark of a civilized people, what would you make of the tradition of cremating the dead, which appeared in many ancient civilizations and especially in Rome, where the nobles and the military heroes in particular were typically cremated as opposed to buried? (Julius Evola's rather polemic view is that burial is ``telluric" and represents a returning to the Earth, whereas cremation indicates an ascent to the heavens.)


\hfill

\texttt{Cologero on 2015-01-23 at 15:25 said: }

Cremation was not practiced in Rome until late into the Republic. Fustel de Coulanges describes the most ancient practices:

\begin{quotex}
From this primitive belief came the necssity of burial. In order that the soul might be confined to this subterranean abode, which was suited to its second life, it was necessary that the body to which it remained attached should be covered with earth. The soul that had no tomb had not dwelling place. It was a wandering spirit. 

\end{quotex}
Re Evola, riddle me this, Tom:

\begin{itemize}
\item He rejected Providence: Masters of the Right: Joseph de Maistre 
\item He rejected marriage and left no heir, rather strange for an alleged baron 
\item He chose cremation instead of burial 
\end{itemize}

\hfill

\texttt{Cologero on 2015-01-23 at 15:33 said: }

White Rabbit, who said sex in itself is a form of mind control? Sex is necessary for the propagation of the race and bonds families and peoples together with a common inheritance. The point made was that alleged ``sexual freedom" leads to the bondage of addiction, and a partner is not necessary for that. I can't comment on your friends.


\hfill

\texttt{ja on 2015-01-23 at 16:49 said: }

Cologero, do you believe Evola was a hypocrite, was not truly traditional, or was simply mistaken on several points ?


\hfill

\texttt{Olavsson on 2015-01-23 at 20:27 said: }

Excellent article, as always.

But regarding the discussion about burial practices… Cremation seems to have been quite a widespread practice in various Indo-European cultures, and still is in Hinduism. So even though cremation is unorthodox seen in light of the European Christian tradition, and might have been untraditional in some European `pagan' cultures as well, I can't see how it would be untraditional in a wider sense. Maybe there is something I have not understood properly, but ancient cremation practices have always seemed to me as manifestations of a rather high religious spirit. To me it represents purification of lower residues by fire – by light; liberation of the spirit from ties slowing its ascension; the impermanence of the body and the ego; and so on. I find its symbolism to be quite beautiful, and have long thought that I will have my own body cremated once I am dead. But I might change my mind if I become convinced that this is actually less favourable than I thought in principle.

Marriage and heir: While Evola wasn't married and didn't reproduce, the same can be said of quite a few historical sages, ascetics, and other great men who weren't really untraditional by any stretch of the imagination. But Evola, though a genius he was, cannot be entirely compared to such figures of course. And to his disadvantage might be pointed out that he probably wasn't celibate either. Evola surely wasn't a perfect man; but I don't think the fact that he chose cremation should be held against him or rejected as untraditional in principle, even if that is true for the particular Catholic tradition he was born into.


\hfill

\texttt{White Rabbit on 2015-01-24 at 00:06 said: }

Cologero: poor phrasing on my part maybe, I didn't intend to imply that. I was noting that ``you need to get laid" is known to be resorted to as a response to political angst.


\hfill

\texttt{August on 2015-01-24 at 06:20 said: }

My opinion on the simplest way to get over questions about how to live is to adopt an orthodox tradition, generally a religion, and follow the established rules sincerely and to the best of one's ability. Develop further according to circumstance and capacity within that framework.

Why worry about whether cremation or burial is better, how each influences the soul, what is permissible sexual behaviour, etc.? These questions are not to be answered absolutely, universally. If you have a tradition, follow the rules and there is your peace.

But many are those who will not, in today's world. If they decide to live outside of a functioning tradition, the only way to ultimately appraise their actions is by the consequences. Given the limits of human horizons, such a path entails placing oneself at the mercy of powers that we cannot grasp, the same powers that know and deliver all men according to some justice, all men naturally including those who performed rites, and those who did not.


\hfill

\texttt{Cologero on 2015-01-24 at 09:58 said: }

Thanks, August, for saving me the trouble. Vico's point is the same whether or not burial or cremation was part of the death rites; he was addressing ancient Europe specifically. And Olavsson, your perspective is sound as I was not trying to define any hard boundaries. The comments just sounded like a series of ``gotcha" attempts while ignoring the main point. Getting a fact wrong here or there is not the issue.

As for the other question, no, I don't think Evola was a hypocrite. Evola, unlike Guenon, chose to engage with the modern world. He was involved politically, culturally, artistically, philosophically and so on. That is why he is more accessible to people today looking for Tradition. On the other hand, the temptation is to assume that familiarity with Evola's writings makes one ``traditional" in some sense. Evola made a conscious choice not to follow an exoteric tradition. This in many ways has distorted things and worst of all it has encouraged others to do the same. Yet he considered the Revolution of 1789 to be the end of civilization, while the rest of contemporary Europe regards it as the beginning of civilization. Who, on any of the various ``rights", will propose blasphemy as a capital crime? And how can Evola prefer pre-1789 Europe while rejecting its spiritual source?

Although Evola regarded Guenon as the ``master of the 20th century", who among the Evola readers bother with Guenon? I only know of one, but please let me know of others. Now Guenon regarded the entire modern world as based on an illusion so he seldom bothered with it. Perhaps occasional book reviews and so on, but nothing sustained in a positive sense. I'd rather see a discussion on this point:

``If our contemporaries as a whole could see what it is that is guiding them and where they are really going, the modern world would at once cease to exist as such"

That is what I've been trying to do and even if it has been done poorly, nevertheless, in Chesterton's phrase, it is still worth doing.


\hfill

\texttt{ja on 2015-01-24 at 13:16 said: }

One aspect of the modern mindset I've noticed on Gornahoor…..Cologero, do you notice how the mention of political topics even as just an aside gathers much commenting whereas spiritual topics are less commented ? The modern man, as per Marx, is seeking to change the world when he should be changing himself. ..


\hfill


\end{sffamily}\end{footnotesize}
